%-------------------------------------------------------------------------
\chapter{Materiais e Métodos}

\section{Caracterização da pesquisa}

Este trabalho caracteriza-se como uma pesquisa aplicada, de natureza tecnológica, voltada ao desenvolvimento e validação funcional de um protótipo de software para apoio ao inventário patrimonial. A pesquisa é aplicada porque busca responder a um problema prático identificado no contexto do Colegiado de Ciência da Computação da UESC: a necessidade de tornar o controle de ativos mais rastreável, organizado e menos dependente de conferência manual.

Quanto aos procedimentos, o trabalho envolve revisão bibliográfica, levantamento de requisitos, modelagem da solução, desenvolvimento de software, integração com dispositivo RFID de baixo custo e validação funcional em cenário controlado. A avaliação proposta concentra-se na verificação do funcionamento do fluxo computacional da solução, não na medição de desempenho de uma infraestrutura RFID de longo alcance. Essa delimitação é coerente com abordagens de verificação e validação de sistemas, nas quais se avalia se o produto atende aos requisitos definidos e ao uso pretendido dentro do escopo estabelecido \cite{ieee1012}.

\section{Materiais utilizados}

A solução foi desenvolvida como uma aplicação web composta por frontend, backend com API REST, banco de dados e processamento de eventos RFID. O frontend foi implementado com tecnologias baseadas em React, Next.js e TypeScript, com o objetivo de fornecer telas para cadastro, consulta, auditoria e acompanhamento operacional. O backend foi desenvolvido em Python com Django e Django REST Framework, concentrando as regras de negócio e expondo endpoints para autenticação, cadastros, eventos RFID, auditoria, logs e inconsistências. A API REST, portanto, integra o backend e atua como ponto padronizado de comunicação entre a interface web, os dispositivos de leitura e os componentes de integração.

Para persistência dos dados, utilizou-se banco SQLite no ambiente de prototipação. Essa escolha é adequada para validação local e simplificação da implantação inicial, embora não represente a configuração ideal para um ambiente institucional em produção. Em uma implantação futura, o banco poderia ser substituído por uma solução mais robusta, como PostgreSQL, sem alterar a proposta conceitual da arquitetura.

Na validação prática, utilizou-se um computador pessoal, um leitor RFID de proximidade de baixo custo disponível para o experimento e uma única tag RFID. Esse leitor opera como dispositivo de entrada local, exige aproximação direta da tag para captura do identificador e não dispõe de interface própria de rede ou mecanismo nativo para envio de requisições HTTP. Por esse motivo, foi utilizado um adaptador intermediário de integração para capturar a tag lida e encaminhar o evento à API REST do sistema. Esse componente foi empregado por uma limitação de interface do equipamento utilizado, e não por restrição da API proposta. A escolha desse equipamento foi orientada por disponibilidade, baixo custo e viabilidade experimental, permitindo validar o pipeline lógico do sistema. A arquitetura, entretanto, foi concebida para aceitar outros dispositivos de captura, desde que estes enviem eventos compatíveis à API. Essa decisão também considera que a implantação física de RFID envolve variáveis específicas de ambiente e interferência, que demandam experimentos próprios para avaliação de desempenho \cite{knapp2022,ting2013}.

\section{Arquitetura proposta}

A arquitetura da solução foi organizada em camadas. A primeira camada corresponde à interface web, responsável pela interação com o usuário. Nela são realizados cadastros de locais, leitores RFID e itens patrimoniais, além da consulta de patrimônio, auditoria, inconsistências e log operacional.

A segunda camada corresponde ao backend com API REST, responsável por receber requisições do frontend e eventos de dispositivos RFID. Essa camada concentra regras de validação, autenticação, ativação de janelas de leitura, processamento de tags, geração de timeline e tratamento de inconsistências. A API é destacada dentro do backend por ser a interface padronizada de entrada e saída de dados do sistema. A terceira camada é o banco de dados, no qual são armazenadas as entidades principais: locais, leitores, itens patrimoniais, leituras, eventos de timeline, auditorias e notificações de inconsistência. A adoção de interfaces padronizadas para integrar dispositivos heterogêneos e aplicações é compatível com requisitos frequentemente discutidos em arquiteturas de middleware para IoT \cite{razzaque2015}.

A quarta camada é a integração com dispositivos de leitura. No cenário escalável previsto, sensores de presença, como sensores infravermelhos, podem acionar a API para abrir uma janela de leitura de determinada antena RFID. Em seguida, as tags capturadas pela antena são enviadas ao backend para processamento. Na validação realizada, o mesmo princípio foi reproduzido com um leitor RFID de proximidade e um adaptador intermediário de integração, responsável por converter a leitura local da tag em uma requisição HTTP compatível com a API REST.

\section{Procedimentos de desenvolvimento}

O desenvolvimento foi conduzido a partir da definição das entidades e fluxos principais do domínio. Inicialmente, foram modelados locais, leitores RFID e itens patrimoniais. Cada item possui uma identificação de etiqueta, um local lógico esperado e um local físico detectado. Essa distinção permite comparar o registro do sistema com a situação observada nas leituras.

Em seguida, foram implementados os endpoints da API e as telas do frontend. O sistema permite cadastrar e consultar itens, acionar leitores, registrar eventos RFID, consultar logs, executar auditorias e listar inconsistências. Também foi implementado um mecanismo para abertura de janelas de leitura, de modo que uma tag recebida seja processada somente quando associada a um leitor ativo ou a uma operação de auditoria.

O processamento das leituras contempla normalização de tags, deduplicação de eventos em intervalo curto, verificação de tags desconhecidas, atualização de localização física e geração de registros históricos. Em leitores configurados como destino, a leitura pode atualizar o local físico do item. Em leitores configurados como fluxo, a leitura pode registrar passagem sem necessariamente alterar o local físico. Essa separação favorece a expansão futura da solução para diferentes pontos de captura.

\section{Metodologia de validação}

A validação foi definida como funcional e qualitativa, pois o objetivo foi verificar se os principais fluxos do sistema operam corretamente com o hardware disponível. Os testes buscaram verificar o fluxo lógico da aplicação: leitura da tag, captura pelo adaptador intermediário, envio do evento à API REST, processamento pelo backend, persistência dos dados e atualização das informações de inventário. Não foram assumidos, sem medição específica, resultados de alcance, leitura simultânea de múltiplas tags, taxa de leitura em massa ou desempenho de antenas RFID de longo alcance. A definição de cenários e resultados esperados aproxima-se de práticas de documentação e execução de testes de software, nas quais casos de teste são estruturados para verificar comportamentos esperados do sistema \cite{iso29119}.

Os cenários de teste previstos incluem: cadastro de locais, leitores e itens; leitura de tag cadastrada; atualização do local físico; geração de histórico operacional; detecção de tag desconhecida; identificação de item em local divergente; tratamento de leituras duplicadas; auditoria com itens encontrados e não encontrados; e comportamento diante de leitor offline. Esses cenários permitem avaliar se o sistema cumpre a função de receber eventos, processá-los e refletir seus efeitos no inventário patrimonial, em alinhamento com a verificação de conformidade entre requisitos, implementação e uso pretendido \cite{ieee1012}.

Quando houver dados quantitativos disponíveis, recomenda-se registrar quantidade de itens testados, quantidade de leituras corretas, leituras duplicadas ignoradas, inconsistências detectadas, duração da janela de leitura e tempo aproximado de resposta da API. Na ausência dessas medições, os resultados devem ser apresentados como validação funcional, com evidências por prints, logs e descrição dos comportamentos observados.
